
\subsubsection{5.1 RQ1: Proxy FAIR Score Variance in the OpenML Corpus}

\begin{figure}[htbp]
\centering
\includegraphics[width=\columnwidth]{figures/fair_distribution.png}
\caption{FAIR score distribution}
\label{fig:fair_distribution}
\end{figure}

\textbf{Figure 1.} Distribution of proxy FAIR compliance scores across 5,000 OpenML datasets. The bimodal distribution (dip test p = 9.96e-6) confirms sufficient variance for matched-pairs analysis. Mean = 0.430, std = 0.069. 85.6\% of datasets score below 0.5.

All existence gate criteria pass:

\begin{table}[htbp]
\centering
\begin{tabular}{llll}
\toprule
Metric & Value & Gate & Status \\
\midrule
CV of proxy FAIR scores & 0.1597 & > 0.15 & **PASS** \\
Bimodality (dip test p) & 9.96e-6 & < 0.05 & **PASS** \\
n_high (score ≥ 0.5) & 720 (14.4\%) & ≥ 500 & **PASS** \\
n_low (score < 0.5) & 4,280 (85.6\%) & ≥ 500 & **PASS** \\
\bottomrule
\end{tabular}
\end{table}

The CV of 0.1597 is a marginal pass (threshold 0.15). The mean score of 0.430 indicates that the corpus is predominantly below the 0.5 threshold, with only 14.4\% of datasets classified as high-FAIR under the proxy measure. Bimodality (Bicoeficient = 0.304, dip test p = 9.96e-6) is statistically confirmed. The weak Spearman correlation between proxy score and OpenML quality metrics (r = 0.055) indicates that the scores are not dominated by dataset size alone.

These results are qualified by the proxy limitation: the normalized quality proxy captures structural/computational characteristics only and does not correspond directly to true F-UJI sub-criteria scores. The near-threshold CV should be re-verified with actual F-UJI scoring.

\subsubsection{5.2 RQ2: Matched vs. Unadjusted Analysis — The Confounding Reversal}

The central empirical finding is the contrast between unadjusted and matched analyses on the same data:

\begin{table}[htbp]
\centering
\begin{tabular}{lll}
\toprule
Analysis & Log-rank p & Interpretation \\
\midrule
Unadjusted KM (full corpus) & 0.583 & Not significant — naive analysis yields null result \\
Matched KM (35 pairs, SMD max = 0.098) & **0.0053** & Significant after confounder adjustment \\
\bottomrule
\end{tabular}
\end{table}

\begin{figure}[htbp]
\centering
\includegraphics[width=\columnwidth]{figures/fig2_km_curves_matched.png}
\caption{Kaplan-Meier survival curves (matched cohort)}
\label{fig:fig2_km_curves_matched}
\end{figure}

\textbf{Figure 2.} Kaplan-Meier survival curves for TTFR in matched pairs (n = 35 pairs). High-FAIR group (solid) reaches median TTFR at 158 days; low-FAIR group (dashed) at 202 days. Matched log-rank p = 0.0053. These results derive from a synthetic smoke-test cohort.

Median TTFR for the high-FAIR group is 158 days, compared to 202 days for the matched low-FAIR group — a difference of 44 days (22\% relative reduction with respect to the low-FAIR median). Covariate balance after matching was confirmed: SMD max = 0.098 < 0.1 threshold for all covariates.

\begin{figure}[htbp]
\centering
\includegraphics[width=\columnwidth]{figures/fig4_love_plot.png}
\caption{Love plot (covariate balance)}
\label{fig:fig4_love_plot}
\end{figure}

\textbf{Figure 3.} Love plot showing standardized mean differences before and after propensity score matching. Post-matching SMD max = 0.098 confirms balance on all three covariates.

Cox proportional hazards results:

\begin{table}[htbp]
\centering
\begin{tabular}{llll}
\toprule
Metric & Value & Gate & Status \\
\midrule
Cox HR & 3.159 & > 1.2 & **PASS** \\
95\% CI & [1.032, 9.672] & Excludes 1.0 & **PASS** \\
Cox p-value & 0.044 & — & Significant \\
\bottomrule
\end{tabular}
\end{table}

\begin{figure}[htbp]
\centering
\includegraphics[width=\columnwidth]{figures/fig5_cox_forest.png}
\caption{Cox forest plot}
\label{fig:fig5_cox_forest}
\end{figure}

\textbf{Figure 4.} Cox proportional hazards forest plot. HR = 3.159 [95\% CI: 1.032, 9.672]. The wide confidence interval reflects the small smoke-test sample (n = 35 matched pairs).

A hazard ratio of 3.159 indicates that high-FAIR datasets have approximately three times the instantaneous rate of receiving a first experimental run compared to matched low-FAIR counterparts. The wide confidence interval [1.032, 9.672] reflects the small sample and underscores the preliminary nature of this estimate. Additionally, the Schoenfeld residuals test flagged a potential PH assumption violation, suggesting that the hazard ratio may not be constant across the observation window; the reported HR should be interpreted as an average effect estimate.

These results derive entirely from a synthetic proof-of-concept cohort (n = 200, 35 matched pairs). The suppressor confounding demonstration — the p-value reversal from 0.583 to 0.0053 on the same data — constitutes the primary methodological contribution. Whether the magnitude of the effect (HR ≈ 3.16) replicates at production scale on real OpenML data with real upload dates and F-UJI scores remains to be determined.

\subsubsection{5.3 RQ3: Findable Sub-criteria Disaggregation vs. Aggregate Scoring}

\begin{table}[htbp]
\centering
\begin{tabular}{llll}
\toprule
Analysis & Log-rank p & Cox HR & Interpretation \\
\midrule
**Findable IV (primary)** & **0.0053** & **3.159** & Significant \\
Ablation A: Aggregate threshold & 0.697 & 1.06 & Near null — not significant \\
Ablation B: Accessible sub-criteria IV & 0.064 & 1.79 & Marginal, not significant at α = 0.05 \\
Ablation C: Relaxed caliper & 0.000 & 3.66 & Consistent with primary \\
\bottomrule
\end{tabular}
\end{table}

Aggregate FAIR threshold scoring (p = 0.697, HR = 1.06) and Findable sub-criteria scoring (p = 0.0053, HR = 3.159) are applied to the same matched pairs. The aggregate scoring result is statistically indistinguishable from the null hypothesis. This suggests that different FAIR dimensions have heterogeneous effects on discovery dynamics, and that aggregate scores dilute the Findable signal. The Accessible sub-criteria ablation (Ablation B: p = 0.064, HR = 1.79) shows a marginal trend that does not reach the pre-specified significance threshold.

\subsubsection{5.4 Sensitivity Analysis: Observation Window}

\begin{figure}[htbp]
\centering
\includegraphics[width=\columnwidth]{figures/fig6_sensitivity_comparison.png}
\caption{Sensitivity analysis panel}
\label{fig:fig6_sensitivity_comparison}
\end{figure}

\textbf{Figure 5.} Sensitivity analysis across observation windows. Log-rank p and Cox HR are consistent across 365-day, 730-day, and 1095-day windows.

\begin{table}[htbp]
\centering
\begin{tabular}{lll}
\toprule
Window & Log-rank p & Cox HR \\
\midrule
365 days & 0.006 & 3.00 \\
**730 days (primary)** & **0.0053** & **3.159** \\
1095 days & 0.005 & 2.93 \\
\bottomrule
\end{tabular}
\end{table}

The Findable effect on TTFR is consistent across 1–3 year observation windows on the smoke-test cohort, ruling out sensitivity to the specific window choice within that range.

\subsubsection{5.5 Summary of Results}

\begin{table}[htbp]
\centering
\begin{tabular}{lll}
\toprule
Research Question & Finding & Confidence \\
\midrule
RQ1: Sufficient FAIR variance exists & CV = 0.1597 (marginal), bimodal (dip p = 9.96e-6), n_high = 720 & Medium — proxy only; near-threshold CV \\
RQ2: Findable → shorter TTFR (matched) & Matched p = 0.0053; HR = 3.159 [1.032, 9.672]; 44-day median reduction & Medium — synthetic cohort, proxy IV, PH violation \\
RQ3: Sub-criteria > aggregate scoring & Findable p = 0.0053 vs. aggregate p = 0.697 & Medium — same caveat on synthetic data \\
Methodological: matching necessity & p = 0.583 → p = 0.0053 on same data & Medium-High — directionally robust \\
\bottomrule
\end{tabular}
\end{table}

\subsubsection{5.6 H-M2: Accessible Dimension — Production Failure}

The Accessible sub-criteria analysis (H-M2) was implemented and validated in dry-run conditions (MWU p = 6.99e-9, standardized β = 0.743 on a synthetic n = 200 cohort with 43 matched pairs, SMD max = 0.092). The production run failed: near-uniform propensity scores (range 0.485–0.515) yielded only 4 matched pairs (target ≥ 500), rendering statistical tests uninformative (MWU p = 1.000, β = −0.042). The root cause is that \texttt{upload_date} was unavailable in the h-e1 fair_scores.csv (all NaN under the proxy scoring pass), so the temporal covariate used in matching lacked real variation. The dry-run success reflects correct pipeline implementation, not hypothesis validity; the true Accessible dimension effect on 12-month run accumulation remains untested at production scale.

